\documentclass[11pt,a4paper]{article}

\usepackage[utf8]{inputenc}
\usepackage[T1]{fontenc}
\usepackage{XCharter}
\usepackage[brazil]{babel}
\usepackage[margin=2cm]{geometry}
\usepackage{amsmath, amssymb}
\usepackage[xcharter,bigdelims,vvarbb]{newtxmath}
% newtxmath's operators font requests the fontspec family `XCharter(0)' in OT1;
% without these substitutions math digits and operator names silently fall back
% to Computer Modern (LaTeX Font Warning: Font shape `OT1/XCharter(0)/m/n' undefined).
\DeclareFontFamily{OT1}{XCharter(0)}{}
\DeclareFontShape{OT1}{XCharter(0)}{m}{n}{<->ssub * XCharter-TLF/m/n}{}
\DeclareFontShape{OT1}{XCharter(0)}{m}{it}{<->ssub * XCharter-TLF/m/it}{}
\DeclareFontShape{OT1}{XCharter(0)}{b}{n}{<->ssub * XCharter-TLF/b/n}{}
\DeclareFontShape{OT1}{XCharter(0)}{bx}{n}{<->ssub * XCharter-TLF/b/n}{}
\usepackage{enumerate}
\usepackage{array}
\usepackage{tikz}
\usetikzlibrary{positioning}

\pagestyle{empty}
\setlength{\parindent}{0pt}
\setlength{\parskip}{0.5em}

\newcommand{\thetav}{\boldsymbol{\theta}}

\renewcommand{\arraystretch}{1.4}

\begin{document}

%% =============================================================
%%  Header
%% =============================================================
{\Large\bfseries Aprendizado Profundo} \hfill \textbf{Prova G2, Unidades 4 e 5}\\
Prof. Lucas S. Kupssinskü \hfill Data: \rule{3cm}{0.4pt}\\[0.4em]
Nome: \rule{12cm}{0.4pt}

\rule{\linewidth}{0.8pt}

%% =============================================================
%%  Response grid
%% =============================================================
\textbf{Quadro de respostas (questões objetivas):}\\[0.3em]
\begin{center}
\begin{tabular}{|c|c|c|c|c|c|c|c|c|}
\hline
\textbf{Questão} & \textbf{a} & \textbf{b} & \textbf{c} & \textbf{d} & \textbf{e} & \textbf{f} & \textbf{g} & \textbf{h} \\\hline
1 & & & & & & & & \\\hline
2 & & & & & & & & \\\hline
3 & & & & & & & & \\\hline
4 & & & & & & & & \\\hline
5 & & & & & & & & \\\hline
6 & & & & & & & & \\\hline
7 & & & & & & & & \\\hline
8 & & & & & & & & \\\hline
\end{tabular}
\end{center}

\rule{\linewidth}{0.4pt}

%% =============================================================
%%  Objective questions  (8 x 1,0 = 8,0 pts)
%% =============================================================
\section*{Questões Objetivas (8{,}0 pts)}
\textit{Cada questão possui apenas uma alternativa correta. Marque a letra escolhida no quadro de respostas.}

\begin{enumerate}[1)]

%% ---------- Q1 (U4): numeric — LoRA trainable parameters ----------
\item (1{,}0) No LoRA, uma matriz congelada $\mathbf{W}_0 \in \mathbb{R}^{d \times k}$ recebe a correção \textit{low rank} $\mathbf{h} = \mathbf{W}_0 \mathbf{x} + \frac{\alpha}{r} \mathbf{B} \mathbf{A} \mathbf{x}$, com $\mathbf{B} \in \mathbb{R}^{d \times r}$ e $\mathbf{A} \in \mathbb{R}^{r \times k}$; apenas $\mathbf{B}$ e $\mathbf{A}$ são treináveis. Você aplica LoRA com $r = 8$ às projeções $\mathbf{W}_Q$ e $\mathbf{W}_V$ (ambas $1024 \times 1024$) de \textbf{cada uma} das $12$ camadas de um \textit{transformer}. Quantos parâmetros treináveis o ajuste introduz no total?

\begin{enumerate}[a)]
    \item $16\,384$
    \item $32\,768$
    \item $98\,304$
    \item $196\,608$
    \item $393\,216$
    \item $589\,824$
    \item $786\,432$
    \item $25\,165\,824$
\end{enumerate}

%% ---------- Q2 (U4): OBJ-B I/II/III — LLM alignment pipeline ----------
\pagebreak
\item (1{,}0) Sobre as asserções abaixo a respeito do treinamento e alinhamento de LLMs:

\textbf{I.} No \textit{SFT} (\textit{supervised fine-tuning}) com \textit{prompt masking}, a função de perda é computada apenas sobre os \textit{tokens} da resposta; os \textit{tokens} do \textit{prompt} são mascarados e não contribuem para o gradiente.

\textbf{II.} O DPO otimiza diretamente sobre pares de preferência $(y_w, y_l)$, dispensando o treinamento de um \textit{reward model} explícito e o loop de \textit{policy gradient} usado no RLHF com PPO.

\textbf{III.} No RLHF com PPO, o termo de penalização $\mathrm{KL}(\pi_\theta \,\|\, \pi_{\text{ref}})$ tem como função principal acelerar a convergência, permitindo o uso de taxas de aprendizado maiores na atualização da política.

\begin{enumerate}[a)]
    \item Apenas I está correta.
    \item Apenas II está correta.
    \item Apenas III está correta.
    \item Apenas I e II estão corretas.
    \item Apenas I e III estão corretas.
    \item Apenas II e III estão corretas.
    \item Todas estão corretas.
    \item Nenhuma está correta.
\end{enumerate}

%% ---------- Q3 (U4): OBJ-C scenario diagnostic + TikZ — catastrophic forgetting ----------
\item (1{,}0) Você faz \textit{fine-tuning} completo (todos os pesos) de um LLM pré-treinado em um corpus jurídico, com taxa de aprendizado $3 \cdot 10^{-4}$ por $10$ épocas. As curvas abaixo mostram a acurácia ao longo das épocas em duas avaliações, ambas sobre dados de teste não vistos no treino: um \textit{benchmark} da tarefa jurídica alvo (linha cheia) e um \textit{benchmark} de conhecimentos gerais que o modelo dominava antes do ajuste (linha tracejada). Qual o diagnóstico mais consistente e a melhor mitigação?

\begin{center}
\begin{tikzpicture}[>=stealth, every node/.style={font=\small}, xscale=0.55, yscale=3.2]
    \draw[->] (0,0) -- (11.2,0) node[below left] {épocas};
    \draw[->] (0,0) -- (0,1.12) node[above] {acurácia};
    \draw (0,1) node[left] {$1{,}0$} -- (0.15,1);
    \draw (0,0.5) node[left] {$0{,}5$} -- (0.15,0.5);
    % target legal task: rises
    \draw[thick] plot[smooth] coordinates {(0,0.52) (2,0.68) (4,0.78) (6,0.85) (8,0.88) (10,0.90)}
        node[right] {\textit{benchmark} jurídico};
    % general benchmark: falls
    \draw[thick, dashed] plot[smooth] coordinates {(0,0.72) (2,0.66) (4,0.56) (6,0.47) (8,0.40) (10,0.35)}
        node[right] {\textit{benchmark} geral};
\end{tikzpicture}
\end{center}

\begin{enumerate}[a)]
    \item O modelo está em \textit{overfitting} no corpus jurídico; treinar por mais épocas no mesmo corpus faria as duas curvas de acurácia voltarem a subir simultaneamente.
    \item Há vazamento de dados entre treino e teste jurídicos; refazer a divisão dos dados corrige tanto a curva da tarefa jurídica quanto a do \textit{benchmark} geral.
    \item Trata-se de esquecimento catastrófico do pré-treino; mitigações incluem taxa de aprendizado menor, menos épocas, misturar dados gerais ao \textit{fine-tuning} ou usar PEFT (LoRA) com os pesos base congelados.
    \item Trata-se de \textit{reward hacking}: o modelo explora falhas da métrica de avaliação; substituir o \textit{benchmark} geral por outro conjunto de tarefas resolve o problema observado.
    \item O \textit{tokenizer} não cobre o vocabulário jurídico e fragmenta os termos técnicos; retreinar o \textit{tokenizer} recupera simultaneamente as duas curvas de acurácia.
    \item A janela de contexto foi excedida pelos documentos jurídicos longos; truncar as entradas do \textit{fine-tuning} evita a degradação observada no \textit{benchmark} geral.
    \item A taxa de aprendizado alta causou explosão de gradiente e \texttt{NaN} nos pesos; nesse caso ambas as curvas deveriam despencar, logo o gráfico é inconsistente com um treino real.
    \item O modelo está subtreinado no regime de Chinchilla ($D^{*} \approx 20N$); aumentar o volume de dados jurídicos melhora as duas curvas por leis de escala.
\end{enumerate}

%% ---------- Q4 (U4): OBJ-H matching — positional encoding methods ----------
\pagebreak
\item (1{,}0) Associe cada método de \textit{positional encoding} (coluna da esquerda) ao seu mecanismo (coluna da direita).

\begin{center}
\begin{tabular}{p{0.34\linewidth} p{0.58\linewidth}}
\textbf{1.} \textit{PE} sinusoidal & \textbf{(i)} tabela de \textit{embeddings} treináveis somada à entrada, limitada ao comprimento máximo $L_{\text{max}}$ visto no treino. \\
\textbf{2.} \textit{PE} absoluto aprendido & \textbf{(ii)} função fixa e determinística da posição, somada à entrada, sem nenhum parâmetro treinável. \\
\textbf{3.} RoPE & \textbf{(iii)} não altera a entrada: subtrai do \textit{score} de atenção uma penalidade proporcional à distância $|i - j|$, com inclinação própria por cabeça. \\
\textbf{4.} ALiBi & \textbf{(iv)} rotaciona os vetores $\mathbf{q}$ e $\mathbf{k}$ dentro da atenção, fazendo o \textit{score} depender apenas do deslocamento relativo $j - i$. \\
\end{tabular}
\end{center}

\begin{enumerate}[a)]
    \item 1-(ii); \quad 2-(i); \quad 3-(iii); \quad 4-(iv)
    \item 1-(i); \quad 2-(ii); \quad 3-(iv); \quad 4-(iii)
    \item 1-(ii); \quad 2-(iv); \quad 3-(i); \quad 4-(iii)
    \item 1-(i); \quad 2-(ii); \quad 3-(iii); \quad 4-(iv)
    \item 1-(iv); \quad 2-(i); \quad 3-(ii); \quad 4-(iii)
    \item 1-(ii); \quad 2-(iii); \quad 3-(iv); \quad 4-(i)
    \item 1-(ii); \quad 2-(i); \quad 3-(iv); \quad 4-(iii)
    \item 1-(iii); \quad 2-(i); \quad 3-(iv); \quad 4-(ii)
\end{enumerate}

%% ---------- Q5 (U5): numeric — saturating vs non-saturating generator loss gradients ----------
\item (1{,}0) Em uma GAN, seja $d = f[g[\mathbf{z},\thetav],\phi]$ o \textit{logit} do discriminador para uma amostra gerada; a \textit{loss} original do gerador contribui com $L = \ln(1 - \sigma(d))$ e a \textit{non-saturating} com $L_{\text{ns}} = -\ln \sigma(d)$. No início do treino o discriminador rejeita as amostras geradas: $\sigma(d) = 0{,}05$. Sabendo que $\sigma'(d) = \sigma(d)\,(1 - \sigma(d))$, quais os módulos de $\left|\partial L / \partial d\right|$ e $\left|\partial L_{\text{ns}} / \partial d\right|$, e qual \textit{loss} fornece sinal útil ao gerador?

\begin{enumerate}[a)]
    \item $0{,}05$ e $0{,}95$; a \textit{loss} original fornece o maior sinal de aprendizado.
    \item $0{,}95$ e $0{,}05$; a \textit{loss} original fornece o maior sinal de aprendizado.
    \item $0{,}05$ e $0{,}05$; as duas \textit{losses} saturam igualmente no início do treino.
    \item $0{,}05$ e $0{,}95$; a \textit{loss} \textit{non-saturating} fornece o maior sinal de aprendizado.
    \item $0{,}95$ e $0{,}05$; a \textit{loss} \textit{non-saturating} fornece o maior sinal de aprendizado.
    \item $0{,}0475$ e $0{,}0475$; as duas \textit{losses} fornecem exatamente o mesmo sinal.
    \item $0{,}95$ e $0{,}95$; as duas \textit{losses} fornecem sinal forte no início do treino.
    \item $0$ e $1$; apenas a \textit{loss} \textit{non-saturating} tem gradiente definido nesse ponto.
\end{enumerate}

%% ---------- Q6 (U5): OBJ-F assertion-reason — classifier-free guidance ----------
\item (1{,}0) Considere as asserções abaixo e a relação proposta entre elas.

\textbf{I.} No treinamento de um modelo difusor com \textit{classifier-free guidance}, a condição $c$ é descartada (substituída por $\varnothing$) em uma pequena fração dos exemplos, tipicamente com probabilidade próxima de $0{,}1$.

\begin{center}\textbf{PORQUE}\end{center}

\textbf{II.} Desse modo uma única rede aprende tanto a predição condicional $\mathbf{g}_t[\mathbf{z}_t, c]$ quanto a incondicional $\mathbf{g}_t[\mathbf{z}_t, \varnothing]$, que são combinadas na amostragem como $\tilde{\mathbf{g}}_t = (1+w)\,\mathbf{g}_t[\mathbf{z}_t, c] - w\,\mathbf{g}_t[\mathbf{z}_t, \varnothing]$, com $w \geq 0$.

Assinale a alternativa correta:
\begin{enumerate}[a)]
    \item As asserções I e II são verdadeiras, e a II é uma justificativa correta da I.
    \item As asserções I e II são verdadeiras, mas a II não é uma justificativa correta da I.
    \item As asserções I e II são verdadeiras, e a I é uma justificativa correta da II, mas não o contrário.
    \item A asserção I é verdadeira, e a II é falsa.
    \item A asserção I é falsa, e a II é verdadeira.
    \item As asserções I e II são falsas.
    \item A asserção I é verdadeira somente quando o peso de guiamento é $w = 0$, e a II é verdadeira.
    \item A asserção I é falsa para modelos de difusão latente, e a II é falsa.
\end{enumerate}

%% ---------- Q7 (U5): OBJ-G ordering — DDPM training iteration ----------
\item (1{,}0) Considere os passos abaixo de \textbf{uma} iteração de treinamento de um modelo difusor (DDPM), conforme o algoritmo visto em aula:

\begin{itemize}
    \item[\textbf{I.}] Amostrar uma instância $\mathbf{x}$ do conjunto de treinamento.
    \item[\textbf{II.}] Amostrar $t \sim \mathrm{Uniforme}\{1, \dots, T\}$ e $\boldsymbol{\epsilon} \sim \mathcal{N}(0, \mathbf{I})$.
    \item[\textbf{III.}] Computar $\mathbf{z}_t = \sqrt{\alpha_t}\,\mathbf{x} + \sqrt{1 - \alpha_t}\,\boldsymbol{\epsilon}$ (kernel de difusão).
    \item[\textbf{IV.}] Computar a perda $\lVert \mathbf{g}_t[\mathbf{z}_t, \thetav_t] - \boldsymbol{\epsilon}\rVert^{2}$.
    \item[\textbf{V.}] Atualizar os parâmetros por descida de gradiente.
\end{itemize}

Assinale a alternativa que apresenta uma ordem válida para os passos:
\begin{enumerate}[a)]
    \item I, III, II, IV, V
    \item II, III, I, IV, V
    \item I, II, IV, III, V
    \item III, I, II, IV, V
    \item I, II, III, V, IV
    \item V, I, II, III, IV
    \item II, I, IV, III, V
    \item I, II, III, IV, V
\end{enumerate}

%% ---------- Q8 (U5): numeric — change of variables (normalizing flow) ----------
\item (1{,}0) Um \textit{normalizing flow} de camada única usa a transformação linear $Y = 3X + 1$ sobre a variável de base $X \sim \mathcal{N}(0, 1)$. Usando a fórmula de mudança de variável e sabendo que $p_X(0) = \tfrac{1}{\sqrt{2\pi}} \approx 0{,}399$, qual o valor aproximado da densidade $p_Y(1)$?

\begin{enumerate}[a)]
    \item $1{,}197$
    \item $0{,}798$
    \item $0{,}399$
    \item $0{,}242$
    \item $0{,}199$
    \item $0{,}133$
    \item $0{,}081$
    \item $0{,}044$
\end{enumerate}

\end{enumerate}

\rule{\linewidth}{0.4pt}

%% =============================================================
%%  Dissertative questions  (2 x 1,0 = 2,0 pts)
%% =============================================================
\pagebreak
\section*{Questões Dissertativas (2{,}0 pts)}

\begin{enumerate}[1)]
\setcounter{enumi}{8}

%% ---------- Q9 (U4): DISS-A — reward model + DPO by hand ----------
\item (1{,}0) No alinhamento de LLMs por preferências, o \textit{reward model} é treinado com a perda de Bradley--Terry
\[
\mathcal{L}_{RM} = -\log \sigma\!\left(r_\phi(x, y_w) - r_\phi(x, y_l)\right),
\]
e o DPO otimiza diretamente
\[
\mathcal{L}_{DPO} = -\log \sigma\!\left(\beta\left[\log\tfrac{\pi_\theta(y_w \mid x)}{\pi_{\text{ref}}(y_w \mid x)} - \log\tfrac{\pi_\theta(y_l \mid x)}{\pi_{\text{ref}}(y_l \mid x)}\right]\right),
\]
em que $y_w$ é a resposta preferida e $y_l$ a preterida. Pode arredondar para $4$ casas decimais.

\begin{enumerate}[a)]
    \item (0{,}5) Para um par com $r_\phi(x, y_w) = 2{,}2$ e $r_\phi(x, y_l) = 0{,}7$, compute a probabilidade de o \textit{reward model} preferir $y_w$ (modelo de Bradley--Terry) e o valor de $\mathcal{L}_{RM}$.
    \vspace{3cm}

    \item (0{,}5) No DPO com $\beta = 0{,}2$, os log-\textit{ratios} valem $\log\tfrac{\pi_\theta(y_w \mid x)}{\pi_{\text{ref}}(y_w \mid x)} = 1{,}5$ e $\log\tfrac{\pi_\theta(y_l \mid x)}{\pi_{\text{ref}}(y_l \mid x)} = -1{,}0$. Compute $\mathcal{L}_{DPO}$ e explique em uma frase o que o DPO dispensa em relação ao \textit{pipeline} completo do RLHF.
    \vspace{3.4cm}
\end{enumerate}

%% ---------- Q10 (U5): DISS-B — diffusion forward + reverse step ----------
\item (1{,}0) Considere um modelo difusor com $T = 2$ e \textit{noise schedule} $\beta = (0{,}4;\ 0{,}25)$. Pode arredondar para $4$ casas decimais.

\begin{enumerate}[a)]
    \item (0{,}5) Compute $\alpha_1$ e $\alpha_2$ e, para a instância $x = 2$ com ruído $\epsilon = 0{,}5$, o valor de $z_2$ obtido pelo kernel de difusão.
    \vspace{3cm}

    \item (0{,}5) Com a estimativa de ruído $g_2[z_2, \thetav_2] = 0{,}6$ e tomando $\sigma_2 = 0$ (amostragem determinística), compute $\hat{z}_1$ pela equação de amostragem reversa. Em seguida, explique em uma frase o papel do termo $\sigma_t\,\epsilon$ que foi ignorado.
    \vspace{3.4cm}
\end{enumerate}

\end{enumerate}

\end{document}
